\subsection{Review: Simplifying Fractions} 
Before we manipulate rational expressions, we must remember how we manipulate fractions. Recall the following definitions:
\begin{definition}
\begin{itemize}
\item Two fractions are \term{equivalent} if they represent the same number.
\item A fraction is in \term{lowest terms} if its numerator and denominator are as small as possible and it includes at most one negative sign.
\item To \term{simplify} a fraction is to write an equivalent in lowest terms.
\end{itemize}
\end{definition}
We can simplify a fraction by ``canceling'' any common factors of the numerator and denominator.
\begin{example}Simplify the following fractions.
\begin{lpic}{}{6}
\foreach \y/\l/\n/\d in {1/a/12/16,2/b/-15/-10,3/c/18/-21} {
	\fractextright{1}{-2*\y+1}{\l. };
	\fractextleft{1}{-2*\y+1}{$\dfrac{\n}{\d}={}$};
}
\end{lpic}
\end{example}
\noi When we deal with larger numbers, this process is easier if we factor our numerator and denominator into primes.
\begin{example}
Simplify the fraction $\dfrac{252}{540}$ by factoring the numerator and denominator into primes.
\begin{lpic}{}{9}

\end{lpic}
\end{example}